% !TeX root = ../main_skripsi.tex
\chapter{PENDAHULUAN}

\section{Latar Belakang Masalah}

[Uraikan kesenjangan antara kondisi empiris dan kondisi ideal. Sertakan data,
fakta, atau bukti lokasi penelitian yang dapat ditelusuri. Jelaskan urgensi,
dampak masalah, dan alasan ilmiah penelitian.]

\section{Identifikasi Masalah}

[Petakan berbagai kemungkinan penyebab masalah berdasarkan bukti pada latar
belakang. Jangan masukkan masalah yang tidak didukung uraian sebelumnya.]

\section{Batasan Masalah}

[Tetapkan lingkup objek, variabel, lokasi, periode, metode, dan batas generalisasi
secara tegas tanpa mengurangi kebermaknaan penelitian.]

\section{Rumusan Masalah}

[Tulis setiap rumusan dalam kalimat tanya yang jelas, spesifik, dan dapat dijawab
oleh data serta teknik analisis yang dipilih.]

\section{Tujuan Penelitian}

[Nyatakan capaian yang selaras satu per satu dengan rumusan masalah.]

\section{Manfaat Penelitian}

[Jelaskan manfaat teoretis dan manfaat praktis yang masuk akal berdasarkan
lingkup penelitian. Jangan menjanjikan dampak yang tidak akan diukur.]

\section{Asumsi Penelitian}

[Bagian opsional. Nyatakan hanya asumsi substantif atau metodologis yang memang
menjadi pijakan penelitian; hapus bagian ini jika tidak digunakan.]
